% 【本文件写什么】impl 实现层：qemu-riscv64-opensbi
%   - 定位：QEMU RISC-V + OpenSBI 板级组合。
%   - crate 路径：os/components/wateros-platform/platform-impl/impl-qemu-riscv64-opensbi/src/
%   - 如何实现对应 api-v0 trait：关键类型、状态机、数据结构
%   - 算法或硬件交互流程（可用伪代码/流程图）
%   - 本 impl 启用的 Cargo [features] 与 arch feature（impl-riscv64 等）
%   - 与其它 impl 变体的差异、切换 feature 时的影响
%   - 性能/内存假设与已知限制
% 【事实来源】
%   - 上述 crate 源码与 Cargo.toml
%   - os/feature-tree.txt 中指向本 impl 的 feature 链

\subsubsection{impl-qemu-riscv64-opensbi}

\code{impl-qemu-riscv64-opensbi}面向QEMU\code{virt}+OpenSBI：\code{BootArgs}把\code{a0}/\code{a1}解释为hart id与DTB物理地址，\code{BootContext}进一步类型化为\code{CPUHartID}与\code{DTBPA}。\code{PlatformTimeImpl::time\_frequency\_hz}在未探测DTB时回退\code{10\_000\_000} Hz（QEMU\code{virt}常见\code{timebase-frequency}）。

控制台经SBI\code{console\_write\_byte}/\code{console\_write}逐字节输出，\code{console\_flush}为no-op。定时器\code{set\_timer}把绝对tick deadline直接交给SBI\code{set\_timer}，与arch\code{read\_time\_tick}刻度一致。复位映射到SBI\code{system\_reset}，支持关机与冷/热重启类型枚举转换。

\begin{rustcode}
impl PlatformTime for QEMURiscv64OpenSBITime {
    fn time_frequency_hz() -> PlatformTimeResult<u64> {
        const QEMU_TIMEBASE_HZ: u64 = 10_000_000;
        Ok(QEMU_TIMEBASE_HZ)
    }
}

pub fn set_timer(time: PlatformTimerDeadline) -> PlatformDeadlineTimerResult<()> {
    if sbi_set_timer(time.0).is_ok() { Ok(()) }
    else { Err(PlatformDeadlineTimerError::Failure) }
}
\end{rustcode}
